\chapter{Estado del arte}
\label{chap:estadoarte}

\lettrine{L}{a} gestión digital de instalaciones deportivas es un mercado en
rápido crecimiento, impulsado por la popularización del pádel y la demanda de
servicios \textit{on-demand}. En este capítulo se analizan las soluciones
existentes, se evalúan las tecnologías candidatas para el desarrollo del sistema
y se justifica la elección tecnológica adoptada.

\section{Soluciones existentes}
\label{sec:soluciones_existentes}

\subsection{Playtomic}

Playtomic es la plataforma líder en España para la reserva de pistas de pádel y
tenis. Permite a los usuarios buscar pistas disponibles en su zona, realizar
reservas en tiempo real y gestionar sus partidos. Desde el punto de vista del
club, ofrece un panel de administración para gestionar disponibilidad, precios y
bonos.

Sus principales fortalezas son la amplia red de clubes integrados y la experiencia
de usuario en la aplicación móvil. Sin embargo, presenta limitaciones relevantes:
es una solución SaaS cerrada sin posibilidad de personalización, cobra una
comisión por reserva que puede resultar elevada para clubes pequeños, y no ofrece
funcionalidades avanzadas de gestión de torneos integradas con el sistema de
reservas.

\subsection{Smashapp}

Smashapp es una plataforma orientada a la gestión de ligas y torneos de pádel.
Su punto fuerte es el módulo de competición, que permite crear torneos, gestionar
inscripciones y publicar resultados. No obstante, sus capacidades de gestión de
reservas son limitadas y está orientado a competiciones federadas más que a la
gestión cotidiana de un club.

\subsection{Sistemas propios de clubes}

Muchos clubes han desarrollado o contratado sistemas a medida, a menudo
implementados como simples formularios web conectados a una base de datos o
incluso como hojas de cálculo compartidas. Estos sistemas presentan problemas
de mantenibilidad, falta de integración entre módulos y dependencia de un
proveedor concreto.

\subsection{Análisis comparativo}

La tabla~\ref{tab:comparativa_soluciones} resume las funcionalidades principales
de las soluciones analizadas. Playtomic y Smashapp son productos cerrados sin
documentación técnica pública ni código accesible, por lo que la comparación
se basa en el uso de sus aplicaciones y en su documentación de cara al
cliente, no en una auditoría de su implementación; debe leerse como una
valoración cualitativa orientativa —suficiente para justificar la elección
de funcionalidades de PadelCenter— y no como una verificación exhaustiva de
las capacidades reales de cada competidor.

\begin{table}[ht]
\centering
\resizebox{\textwidth}{!}{%
\begin{tabular}{lcccc}
\hline
\textbf{Funcionalidad} & \textbf{Playtomic} & \textbf{Smashapp} & \textbf{Sist.\ propios} & \textbf{PadelCenter} \\
\hline
Reserva de pistas        & Sí       & Limitado  & Variable & Sí \\
Gestión de torneos       & Básico   & Sí        & Raro     & Sí \\
Emparejamientos auto.    & No       & Sí        & No       & Sí \\
Personalizable           & No       & Limitado  & Sí       & Sí \\
Código abierto           & No       & No        & Variable & Sí \\
Sin comisión por reserva & No       & No        & Sí       & Sí \\
Arquitectura moderna     & Descono. & Descono.  & No       & Sí \\
\hline
\end{tabular}}
\caption{Comparativa de soluciones existentes.}
\label{tab:comparativa_soluciones}
\end{table}

\section{Tecnologías evaluadas}
\label{sec:tecnologias_evaluadas}

\subsection{Frameworks de backend}

Se evaluaron tres frameworks Java para el backend:

\paragraph{Spring Boot~\cite{springboot}} Es el framework más maduro y ampliamente adoptado del
ecosistema Java. Ofrece un ecosistema de módulos muy completo (Spring Security,
Spring Data JPA, Spring MVC), una comunidad enorme y excelente integración con
herramientas de generación de código a partir de OpenAPI. Su modelo de
programación es familiar para cualquier desarrollador Java.

\paragraph{Quarkus} Framework diseñado para entornos cloud-native, con tiempos
de arranque muy reducidos y menor huella de memoria. Aunque es una opción muy
interesante para microservicios, su ecosistema es menos maduro que Spring Boot
y la curva de aprendizaje es mayor para quienes provienen del mundo Java EE.

\paragraph{Micronaut} Similar a Quarkus en su orientación cloud-native, con
inyección de dependencias en tiempo de compilación. Sus ventajas en rendimiento
son similares a las de Quarkus, pero su adopción en el mercado es significativamente
menor, lo que implica menor disponibilidad de recursos y librerías.

La elección recayó sobre \textbf{Spring Boot} por su madurez, la calidad del
plugin OpenAPI Generator para Maven y la cobertura de todas las necesidades
del proyecto sin compromisos.

\subsection{Frameworks de frontend}

Se evaluaron tres frameworks JavaScript/TypeScript para el frontend:

\paragraph{Next.js~\cite{nextjs}} Framework React con renderizado híbrido (SSR, SSG y RSC),
App Router en su versión 14, y un ecosistema muy activo. Su modelo de enrutamiento
basado en el sistema de ficheros es intuitivo y facilita la organización del
código. La integración con React Query y la generación de código con Orval es
directa.

\paragraph{Nuxt} Equivalente de Next.js para el ecosistema Vue.js. Ofrece
características similares pero la comunidad de Vue es más pequeña y el autor
tiene mayor experiencia con React.

\paragraph{Angular} Framework completo de Google, con un enfoque más opinionado
y una mayor cantidad de boilerplate. Su curva de aprendizaje es más pronunciada
y su arquitectura, aunque sólida, resulta excesiva para un proyecto de este
tamaño.

La elección recayó sobre \textbf{Next.js} por su flexibilidad, la calidad del
ecosistema React y la compatibilidad con las herramientas de generación de código
elegidas.

\section{Arquitectura hexagonal vs. arquitectura en capas}
\label{sec:arquitectura_comparativa}

La arquitectura en capas tradicional (\textit{Presentation → Service → Repository})
es el patrón más extendido en aplicaciones Spring Boot. Sin embargo, presenta
un problema fundamental: las capas superiores acaban dependiendo de detalles de
infraestructura (JPA, SQL, HTTP) que se filtran a través de las capas, acoplando
el dominio de negocio a tecnologías concretas. En la práctica, el \texttt{Service}
suele inyectar directamente una interfaz \texttt{JpaRepository} de Spring Data:
el dominio de negocio queda comprometido, en tiempo de compilación, con un
\textit{framework} de persistencia concreto.

\subsection{Ports \& Adapters: origen y conceptos}

La arquitectura hexagonal, también conocida como \textit{Ports \& Adapters},
fue propuesta por Alistair Cockburn~\cite{cockburn2005hexagonal} para resolver
precisamente ese problema, partiendo de una observación: tanto la interfaz de
usuario como la base de datos son, desde el punto de vista del núcleo de la
aplicación, igualmente \textit{externas}. La forma hexagonal no tiene un
significado literal de seis lados; representa simplemente que el núcleo puede
exponer tantos lados (puertos) como necesite, a diferencia del rectángulo de
arriba a abajo de la arquitectura en capas, que sugiere una única dirección de
dependencia lineal.

Cockburn distingue dos tipos de puerto:
\begin{description}
  \item[Puertos primarios (\textit{driving}):] interfaces que el núcleo
    \emph{ofrece} hacia el exterior. Un caso de uso (\texttt{CreateBookingUseCase})
    es un puerto primario: algo de fuera —un controlador REST, una prueba
    automatizada, una futura interfaz de línea de comandos— lo invoca para que
    el núcleo haga su trabajo.
  \item[Puertos secundarios (\textit{driven}):] interfaces que el núcleo
    \emph{necesita} del exterior para funcionar. \texttt{BookingRepository} es
    un puerto secundario: el núcleo lo invoca, pero quien lo \emph{implementa}
    está fuera del núcleo.
\end{description}

Cada puerto tiene uno o varios \textit{adaptadores} que lo conectan con una
tecnología concreta. Los adaptadores primarios \emph{dirigen} el núcleo
(un controlador Spring MVC que traduce una petición HTTP en una llamada a
\texttt{CreateBookingUseCase.execute()}); los adaptadores secundarios son
\emph{dirigidos por} el núcleo (\texttt{BookingRepositoryImpl}, que traduce las
llamadas del puerto a consultas JPA). La idea distintiva del patrón —más allá
de la simple inyección de dependencias— es esta \textbf{simetría}: desde la
perspectiva del núcleo da igual si quien lo invoca es un controlador REST, una
cola de mensajes o un test unitario (son adaptadores primarios intercambiables),
y da igual si la persistencia es PostgreSQL, otra base de datos o un doble de
prueba en memoria (son adaptadores secundarios intercambiables). El núcleo
solo conoce las interfaces; nunca conoce qué adaptador concreto está al otro
lado.

Esta misma idea de inversión de dependencias hacia el núcleo es el principio
central de la \textit{Clean Architecture} de Robert C.~Martin~\cite{martin2017cleanarchitecture},
que reformula los dos tipos de puerto de Cockburn como anillos concéntricos
(entidades, casos de uso, adaptadores de interfaz, \textit{frameworks} y
controladores) gobernados por la \textit{Dependency Rule}: las dependencias de
código solo pueden apuntar hacia el interior, nunca hacia afuera. PadelCenter
sigue esta idea de forma directa, organizando el código en tres paquetes que se
corresponden con los conceptos de Cockburn: \texttt{domain} y
\texttt{application} forman el núcleo hexagonal (entidades y casos de uso, que
definen los puertos como interfaces Java); \texttt{infrastructure.inbound}
contiene los adaptadores primarios (controladores REST); e
\texttt{infrastructure.outbound} contiene los adaptadores secundarios
(repositorios JPA, el adaptador de seguridad JWT). El propio patrón de mapeo en
dos capas descrito en la sección~\ref{sec:implementacion_backend} recoge ideas
clásicas de mapeo objeto-relacional descritas por Fowler~\cite{fowler2002peaa},
adaptadas para mantener las entidades JPA completamente fuera del dominio.

\subsection{Un ejemplo concreto: \texttt{BookingRepository}}

El puerto secundario \texttt{BookingRepository} (interfaz en
\texttt{domain.repository}) declara, entre otras, la operación que detecta
conflictos de horario:

\begin{lstlisting}[language=Java]
public interface BookingRepository {
    boolean existsOverlappingBooking(
        UUID fieldId, LocalDateTime start,
        LocalDateTime end, Long excludeId);
    // ...
}
\end{lstlisting}

\texttt{CreateBookingUseCase} depende únicamente de esta interfaz: ni la
importa, ni puede importar, ningún tipo de Spring Data JPA. El adaptador
secundario que la implementa, \texttt{BookingRepositoryImpl}, vive en
\texttt{infrastructure.outbound} y es el único punto del sistema que conoce a
la vez el puerto del dominio y la API de Spring Data JPA:

\begin{lstlisting}[language=Java]
@Repository
public class BookingRepositoryImpl implements BookingRepository {
    private final BookingJpaRepository bookingJpaRepository;
    private final BookingPersistenceMapper mapper;
    // delega en bookingJpaRepository y traduce
    // entidades JPA <-> records del dominio
}
\end{lstlisting}

En una arquitectura en capas tradicional, \texttt{CreateBookingUseCase} (o el
equivalente \texttt{BookingService}) inyectaría directamente
\texttt{BookingJpaRepository}, una interfaz de Spring Data: el dominio
importaría \texttt{org.springframework\-.data\-.jpa\-.repository\-.JpaRepository} y
quedaría acoplado a JPA en tiempo de compilación, exactamente el problema
señalado al inicio de esta sección. Con el puerto de por medio, esa dependencia
desaparece del dominio: Spring conecta caso de uso y adaptador en tiempo de
\emph{ejecución} mediante inyección de dependencias, pero en tiempo de
\emph{compilación} \texttt{application} no depende de \texttt{infrastructure}
en ningún sentido —la regla que ArchUnit verifica automáticamente en cada
build (sección~\ref{sec:tests_arquitectura}).

Las ventajas concretas para PadelCenter, consecuencia directa de esta
inversión, son:
\begin{itemize}
  \item Posibilidad de testar los casos de uso sin levantar el contexto Spring:
        un test sustituye \texttt{BookingRepository} por un doble de prueba
        (Mockito), sin que el caso de uso note la diferencia.
  \item Facilidad para sustituir la base de datos o el mecanismo de autenticación
        sin modificar el dominio: bastaría con escribir un nuevo adaptador que
        implemente el mismo puerto.
  \item Reglas de arquitectura verificables automáticamente con ArchUnit.
\end{itemize}

La desventaja principal es el mayor número de clases e interfaces necesarias
respecto a una arquitectura en capas plana. Sin embargo, en un proyecto con
ambición de mantenibilidad a largo plazo, esta inversión es justificada.
